% Section 12.3, from lectures/L12/README.md.

\section{Summary}
\label{c12:sec:review}

\needspace{6\baselineskip}
\subsection*{What you should be able to do now}
\begin{itemize}
  \item Implement a width-generic two-flip-flop synchroniser, and explain why it needs no reset and
    what widening it does \emph{not} solve.
  \item Implement a counter-based bit timer producing correctly timed \code{sample} and
    \code{bit_done} pulses, with a \code{resync} input.
  \item Explain what \code{resync} does and why \code{sample} comes before \code{bit_done}.
  \item Read a tick-level trace over one bit period and say for every tick what every output should
    be.
  \item Run a testbench handed out all the way through and trace a failure back to the tick in the
    design that caused it.
\end{itemize}

\needspace{6\baselineskip}
\subsection*{Questions to test yourself with}
\begin{enumerate}
  \item Why does a synchroniser need two flip-flops rather than one, and why does it need no reset?
  \item A widened \code{meta_prev} removes metastability but not skew between the bits. What is the
    difference, and when does a synchroniser stop being the right tool?
  \item What does \code{bit_timer.resync} do, and when is it used?
  \item Why does \code{sample} come before \code{bit_done} rather than the other way round?
  \item \code{bit_timer} holds its counter while \code{enable = '0'}. Which state in
    \code{can_controller} is the only one that would ever want that, and why does it turn out not to
    want it either?
  \item When a testbench handed out fails, what does GHDL print, and how do you find which check
    failed?
\end{enumerate}

\needspace{6\baselineskip}
\subsection*{Target}
\code{meta_prev} and \code{bit_timer} should be finished and merged within about a week. Both are
short, and they are the only modules no other module can substitute for. These are also the first
real pull requests: use the review, and check the port order, that the testbench passes, and that no
constants are hardcoded past \code{can_def}.

\needspace{6\baselineskip}
\subsection*{Looking ahead}
\Chapref{c13:ch} builds \code{crc15}, a bit-serial checksum engine that both generates and checks a
CRC with the same logic and with no mode switch.
